% titel: Blatt 1, A3: SI-Einheiten und Dimensionsanalyse
% skript: Kap. Methoden der Physik, Einheiten

% Einheiten von Konstanten, Dimensionsprüfung, Angstroem-Umrechnung.

Die Dimensionsanalyse prüft Gleichungen auf Konsistenz, ohne Zahlen
einzusetzen -- jede physikalische Gleichung muss beidseitig gleiche
Dimension besitzen.

\begin{werkzeugkasten}[Werkzeuge]
\textbf{Basisgroessen und Einheiten} im SI:
L\"ange \unit{\metre}, Masse \unit{\kilogram}, Zeit \unit{\second},
Stromst\"arke \unit{\ampere}, Temperatur \unit{\kelvin}.
Wichtig: \textbf{Newton}
\begin{equation}
  \unit{\newton}=\frac{\unit{\kilogram\metre}}{\unit{\second\squared}},
  \qquad
  \unit{\pascal}=\frac{\unit{\newton}}{\unit{\metre\squared}}
  =\frac{\unit{\kilogram}}{\unit{\metre\second\squared}}.
\end{equation}

\textbf{Regeln:}
\begin{itemize}
  \item In einer Summe/Gleichung m\"ussen alle Terme gleiche Dimension
        haben.
  \item Das Argument von Funktionen wie $\ln$, $\sin$, $\cos$ muss
        dimensionslos sein (hier: $C_2t$ dimensionslos).
  \item Eine Konstante $C$, die multipliziert wird, übernimmt die
        Einheit, die die Gleichung ausgleicht.
\end{itemize}

\textbf{Dezimalvorsaetze} (Potenzen von 10):
$1\,\mathrm{\mu m}=10^{-6}\,\mathrm m$, $1\,\mathrm{nm}=10^{-9}\,\mathrm m$,
$1\,\mathrm{fm}=10^{-15}\,\mathrm m$.
\end{werkzeugkasten}

\begin{werkzeugkasten}[Vorgehen]
\begin{itemize}
  \item Einheitenprodukte links und rechts zerlegen und vergleichen;
        die Konstante bekommt genau die fehlende Einheit.
  \item Fl\"achen/Beträge: Groessenproben auf die gleiche Kette von
        Masse, Länge, Zeit reduzieren.
  \item Für Umrechnungen den Wert als Produkt ``Zahl mal Zehnerpotenz''
        mit Vorsatz schreiben und in \unit{\metre} umrechnen.
\end{itemize}
\end{werkzeugkasten}

\begin{achtung}[Typische Fallen]
\begin{itemize}
  \item $\ln(C_2t)$: nicht die Einheit von $t$ in $C_2$ ``stecken'',
        sondern $C_2$ so wählen, dass $C_2t$ dimensionslos ist.
  \item ``Gravitationskonstante'' ist die Einheit von $C$ in einer
        Kraftgleichung vom Typ $F=C\,m_1m_2/x^2$, nicht $\vec F=m\vec a$.
  \item Vorsaetze beim Umrechnen mit 10er-Potenzen: einer der
        Umrechnungswerte ist inkorrekt; einfach alle durchexerzieren.
\end{itemize}
\end{achtung}
